\section{Irreducible Hsieh support and its exact local-type limit}
\label{sec:cubic-hsieh-stable-component}

Throughout this section the ordinary cuspidal Hecke algebra is the
one actually constructed in Hsieh, \emph{Eisenstein congruence on
unitary groups and Iwasawa main conjectures for CM fields},
version of 16 September 2025.  We retain the hypotheses and published
inputs of Theorem~7.18 and the prescribed auxiliary-level construction
of AC-K29-Hsieh-aux.  In particular, this is a split coefficient-prime
construction, not an implication of Theorem~A.  Write
\[
 \Lambda=\mathcal O[[\Gamma^+\times\Gamma^-]]
          =\mathcal O[[X,Y]],\qquad R=h_E^{\mathrm{ord}},
 \qquad I=I^S(\Psi),                                      \tag{G64.1}
\]
for an imaginary quadratic $K$, where $p>3$ splits in $K$,
$p\nmid h_K$, and the finite branch character $\beta$ has order
prime to $p$ and is unramified above $p$.
The coefficient ring $\mathcal O$ is a complete discrete valuation
ring.  The nonzero Katz and imprimitive Deligne--Ribet functions,
with the integral normalizations of Hsieh \S6.5, are denoted by
$L_K^S\in\Lambda$ and $L_D^S\in\mathcal O[[X]]$.
Let $q$ be the extra split prime in AC-K29-Hsieh-aux, so that for
$w\mid q$ the finite character satisfies
$\beta(\operatorname{Frob}_w)=1$.

\begin{lemma}[AC-G64-prime: the added split Euler factor has a nonexceptional divisor]
\label{AC-G64-prime}
The anticyclotomic projection of $\operatorname{Frob}_w$ in
$\Gamma^-$ is nonzero.  Consequently there is a height-one prime
$P\subset\Lambda$ such that
\[
       \operatorname{ord}_P(L_K^S)>0,
       \qquad \operatorname{ord}_P(L_D^S)=0.                \tag{G64.2}
\]
In particular $P$ is nonexceptional in the precise sense of
Hsieh Definition~8.13.  There is no estimate here for the auxiliary
prime $q$.
\end{lemma}
\begin{proof}
\proofstep{The split-prime Frobenius has a nonzero anticyclotomic direction.}
Let $h=h_K$, and choose $\pi\in K^\times$ with
$w^h=(\pi)$.  Write $p\mathcal O_K=\mathfrak p\bar{\mathfrak p}$.
The ray-class exact sequence for the moduli
$\mathfrak p^n\bar{\mathfrak p}^{\,n}$, followed by inverse limits and pro-$p$
completion, identifies the two local principal-unit directions
with the two directions of the maximal abelian pro-$p$ extension
unramified outside $p$.  Here the global units are finite and
$p\nmid h$, so neither contributes a pro-$p$ quotient or kernel.
Under conjugation, the anticyclotomic direction is the difference
of the two local logarithms.  Applied to the principal ideal
$w^h$, reciprocity therefore identifies the logarithm of its
anticyclotomic Artin symbol, up to a nonzero fixed normalization
and sign, with
\[
                \log_p(\iota_p(\pi/\bar\pi)),              \tag{G64.3}
\]
where $K_{\mathfrak p}=\mathbb Q_p$ and $\iota_p:K\hookrightarrow\mathbb Q_p$
is the chosen embedding.  This description also follows directly
at finite ray level: the principal generator contributes its two
$p$-adic unit residues, and the anti-invariant quotient takes
their ratio.

The argument of the logarithm is a $p$-adic unit.  The kernel of
$\log_p$ on $\mathbb Z_p^\times$ is the finite group of
Teichm\"uller roots of unity: on $1+p\mathbb Z_p$ the logarithm
and exponential are inverse convergent power series, and the
remaining factor is $\mu_{p-1}$.  If (G64.3) vanished, then
$\pi/\bar\pi$ would have finite order, since $\iota_p$ is
injective.  Its principal ideal is $w^h/\bar w^{\,h}$,
which is nontrivial because $w\ne\bar w$.  This contradiction
proves the assertion.  In particular the anticyclotomic exponent
of $\operatorname{Frob}_w$ is a nonzero element of $\mathbb Z_p$.

\proofstep{No cyclotomic-only factor divides this Euler factor.}
Write the group-like element of that Frobenius as
\[
       [\operatorname{Frob}_w]=(1+X)^a(1+Y)^b,
          \qquad a,b\in\mathbb Z_p,\quad b\ne0.
\]
The Euler factor occurring in Hsieh \S6.5, immediately before
(6.14), is
\[
       f_w=1-(1+X)^a(1+Y)^b.                              \tag{G64.4}
\]
It is a nonzero nonunit in the regular local unique factorization
domain $\mathcal O[[X,Y]]$.

Let $d(X)$ be any prime factor in $\mathcal O[[X]]$ of a nonzero
element of that ring.  By Weierstrass preparation it is either a
uniformizer or, up to a unit, a distinguished irreducible polynomial.
The ideal it generates remains prime in $\mathcal O[[X,Y]]$:
the quotient is $(\mathcal O[[X]]/(d))[[Y]]$, a domain.
If $d$ is not a uniformizer, that quotient has characteristic zero.
The coefficient of $Y$ in the image of $f_w$ is
$-(1+X)^a b\ne0$, so $d$ does not divide $f_w$.
If $d$ is a uniformizer, write $b=p^s u$ with
$u\in\mathbb Z_p^\times$.  In characteristic $p$ the first
nonconstant term of $(1+Y)^b$ is
$\bar uY^{p^s}$, so the image of $f_w$ is again nonzero.
Thus $f_w$ has no height-one divisor in common with any nonzero
cyclotomic-only element, in particular with $L_D^S$.

Choose a prime factor of the nonzero nonunit $f_w$ and let $P$
be the corresponding height-one prime.  The integral imprimitive
factorization has $f_w\mid L_K^S$.  The preceding argument gives
$P\nmid L_D^S$, proving (G64.2).
\end{proof}

\begin{proposition}[AC-G64-stable: a component through that congruence is generically absolutely irreducible]
\label{AC-G64-stable}
Assume also Hsieh's residual multiplicity-free condition
\textup{(Dist)}.  For $P$ as in (G64.2), there is a prime
$\mathfrak q\subset R$ with $I\subset\mathfrak q$ and
$\mathfrak q\cap\Lambda=P$.
Every minimal prime $\mathfrak a\subset\mathfrak q$ is a
component on which the generic three-dimensional pseudocharacter
is absolutely irreducible.  That component contains classical
ordinary cuspidal points with absolutely irreducible attached
Galois representation.

This proves existence of irreducible classical support at the
enlarged Hsieh level.  It does not assert the prescribed away-inertia
or symmetric local Hodge type of the source seed, nor that an
arbitrary actual characteristic-$p$ curve kills the Hecke ideal.
\end{proposition}
\begin{proof}
\proofstep{Work at a prime where the congruence ideal is actually proper.}
Hsieh Definition~7.19 and the equality immediately after it give
\[
 R/I\simeq\Lambda/\operatorname{Eis}(\Psi,S),
 \qquad \operatorname{Eis}(\Psi,S)\subset(L_D^SL_K^S)
                                                               \tag{G64.5}
\]
(Corollary~7.20).  Thus $\operatorname{Eis}(\Psi,S)\subset P$,
and the inverse image of $P/\operatorname{Eis}(\Psi,S)$ in $R$
is the desired prime $\mathfrak q$.  Fix a minimal prime
$\mathfrak a\subset\mathfrak q$.
The reduced ring $R$ is finite and torsion-free over $\Lambda$.
Every minimal prime contracts to zero: otherwise a nonzero
element of $\Lambda$ would lie in a minimal prime of $R$, hence
would be a zero divisor of $R$, contradicting torsion-freeness.
Consequently $F_{\mathfrak a}=\operatorname{Frac}(R/\mathfrak a)$
is a field over $\operatorname{Frac}(\Lambda)$.

\proofstep{Use the faithful trace quotient, not an arbitrary triangular lift.}
Hsieh \S8.3 uses the generalized matrix algebra supplied by
Bella\"iche--Chenevier, Theorem~1.4.4(ii), on
\[
                  \mathscr S=R[G_{K,S}]/\operatorname{Ker}T.
                                                               \tag{G64.6}
\]
Here $\operatorname{Ker}T$ is the trace radical: an element $x$
is killed exactly when $T(xy)=0$ for every $y$.
Its matrix entries are finite fractional ideals $R_{ij}$,
and its diagonal idempotents are the three usual matrix units.
The trace pairing remains nondegenerate after passing to the
total fraction ring and then to the factor $F_{\mathfrak a}$.
Indeed, embed an element of that factor as zero in every other
factor.  If it paired to zero with the whole algebra, clear one
common regular denominator to obtain an element of
$\operatorname{Ker}T$ in $\mathscr S$, which is zero by (G64.6).

It follows in particular that if $R_{ij}(\mathfrak a)\ne0$,
then $R_{ji}(\mathfrak a)\ne0$.  To verify this directly, the
generic algebra contains $E_{ij}$ after multiplying a nonzero
entry by a scalar.  Nondegeneracy gives $y$ with
$\operatorname{tr}(E_{ij}y)=y_{ji}\ne0$; multiplying on either
side by the diagonal idempotents isolates $y_{ji}E_{ji}$.

\proofstep{Exclude vanishing of the middle-to-endpoint entry.}
In the notation of Hsieh (8.6),
\[
 I_1=R_{13},\qquad J_1=R_{12}R_{23}+IR_{13},
 \qquad R_{12}=R_{23}.                                   \tag{G64.7}
\]
The last equality is the conjugate-dual symmetry (8.3)--(8.4).
The complete argument immediately before the faithfulness step
in the proof of Hsieh Theorem~8.14(ii), using Corollary~8.9,
Lemma~8.10 and Theorem~8.11, gives
\[
                         (I_1)_P=(J_1)_P                  \tag{G64.8}
\]
under exactly (G64.2).  These are retained published local
Selmer and totally real main-conjecture inputs; no implication
from Theorem~A is being substituted for them.

Suppose $R_{23}(\mathfrak a)=0$.  Symmetry gives
$R_{12}(\mathfrak a)=0$.  Map (G64.8) to the component and
localize at $\mathfrak q/\mathfrak a$.  For the finite module
obtained from $R_{13}(\mathfrak a)$, it becomes $M=IM$.
The ideal $I$ is contained in the maximal ideal of this local
ring, by the choice of $\mathfrak q$.  Nakayama's lemma yields
$M=0$.  Since $R_{13}(\mathfrak a)$ is a fractional ideal in
the domain's fraction field, localization cannot kill a nonzero
element; thus $R_{13}(\mathfrak a)=0$.
Nondegeneracy of the trace pairing also kills the opposite
entries $R_{21}(\mathfrak a)$, $R_{32}(\mathfrak a)$ and
$R_{31}(\mathfrak a)$.
The generic algebra is therefore diagonal, so its trace is a
sum of three characters.

This is excluded by the final automorphic argument in Hsieh's
Theorem~8.14 proof, printed p.~101.  For precision, diagonal
entries now multiply as characters already over $R/\mathfrak a$,
since their off-diagonal products vanish there.  At a classical
point on this component the semisimple attached representation
is therefore a sum of three characters.  Choose a sufficiently
regular classical point, as allowed by the dense classical cone;
the cited non-CAP and discrete-packet results used in Hsieh's
proof then make it endoscopic of type $(1,1,1)$, with each
character individually conjugate self-dual.  Reduction would
make each of $\bar\delta_1,\bar\delta_2,\bar\delta_3$
individually conjugate self-dual.  But conjugate duality
interchanges $\bar\delta_1$ and $\bar\delta_3$, which are
distinct by \textup{(Dist)}.  This contradiction proves
$R_{23}(\mathfrak a)\ne0$.
The non-CAP and packet classification inputs are used here in
their published source scope, not reproved as elementary facts.

\proofstep{The generic image and some classical images are full matrix algebras.}
The preceding nonvanishing and symmetry put $E_{23}$ and
$E_{12}$ in the generic image.  Trace nondegeneracy puts
$E_{32}$ and $E_{21}$ there.  Products give
$E_{13}=E_{12}E_{23}$ and $E_{31}=E_{32}E_{21}$.
Together with the three diagonal idempotents these are all
matrix units, so the generic image is
$M_3(F_{\mathfrak a})$.

Choose nine group elements $g_1,\ldots,g_9$ whose images span
this matrix algebra; this is possible because the image algebra
is the linear span of the group image.  The trace pairing on
$M_3$ is nondegenerate, so
\[
           \Delta=\det\bigl(T(g_i g_j)\bigr)_{1\le i,j\le9}
                   \notin\mathfrak a.                    \tag{G64.9}
\]
Classical ordinary points are Zariski dense on the component
(Hsieh Proposition~7.17 and \S8.3), so some such point has
$\Delta(x)\ne0$.  At that point the nine image matrices have
nondegenerate Gram matrix, hence span $M_3$ over the coefficient
field.  Every invariant subspace is then stable under every
matrix, proving absolute irreducibility of the attached
representation.  This also proves openness of the asserted
irreducibility here without an additional reducibility-locus
theorem.
\end{proof}

\begin{proposition}[AC-G64-weight: those classical points do not have the required symmetric seed type]
\label{AC-G64-weight}
Suppose that the Hsieh residual branch has been identified,
in either permitted endpoint order, with the actual cubic
tuple $\bar\chi^{\,j}\bar\gamma\oplus1\oplus
\bar\chi^{\,j}\bar\gamma^{-1}$, and that the three residual
inertia characters above $p$ are distinct.  The ordinary
classical points in the separable cone of Hsieh \S8.4 cannot
be transported by a scalar de Rham character and a permutation
to the source seed's type
\[
       \text{residual middle }1\text{ of weight }0,
             \qquad\text{endpoint weights }m,-m,
             \qquad m>0,                                 \tag{G64.10}
\]
while preserving that residual tuple.  The same conclusion
holds after a finite base change on which $\bar\chi$ remains
nontrivial and these residual inertia characters remain distinct.
It does not exclude a different
companion-form, refinement, or local-type theorem.
\end{proposition}
\begin{proof}
At $w\in\Sigma_p$, Hsieh \S8.4, printed p.~98, gives
the ordinary graded characters in the order
$(\delta'_3,\delta'_1,\delta'_2)$, with arithmetic
Hodge--Tate weights
\[
              (k+a-1,\ a+1,\ 0),\qquad k>2,\ a\ge0.
                                                               \tag{G64.11}
\]
Their residual labels are $(\bar\delta_3,\bar\delta_1,
\bar\delta_2)$, where $\bar\delta_2=1$.
In particular the unique graded character with residual
inertia $1$ has weight zero, and the other two weights are
strictly positive.

Let $\theta$ be a scalar twist preserving the displayed
global residual multiset.  Taking its residual determinant
gives $\bar\theta^3=1$.  If $\bar\theta\ne1$, multiplication
by it permutes the three residual characters freely; since
one character is $1$, their product would be $1$.  Their
actual product is $\bar\chi^{\,2j}\ne1$.  Thus
$\bar\theta=1$, exactly as in the residual part of
AC-G63-selftwist.

Let $b$ be the Hodge--Tate weight of $\theta$ at the chosen
embedding.  After twisting, the unique residual-middle
graded character has weight $b$.  Any complete local flag
has the same one-dimensional Jordan--H\"older factors;
distinct residual inertia labels prevent a permutation from
assigning this factor to an endpoint.  Condition (G64.10)
therefore forces $b=0$.  The two remaining weights are
still $a+1>0$ and $k+a-1>0$, and cannot be $m,-m$.
Changing the global Hodge--Tate sign convention negates
all three weights and gives the same contradiction.
Restriction to a finite local extension preserves the
Hodge--Tate weights at every extended embedding.  The retained
nontrivial cubic character supplies the same determinant
contradiction for a scalar twist over the enlarged field, and
the retained residual inertia distinction preserves the label
argument.  These are precisely the two extra hypotheses in
the base-change assertion.
\end{proof}

\paragraph{Exact remaining arithmetic interface.}
The primary inputs used above yield an irreducible component
at the enlarged level, and AC-G63-selftwist preserves its
irreducibility under finite abelian base change.  At the added
prime $q$, however, the actual cubic residual datum is
unramified.  The literal killing condition of the actual
deformation problem consequently requires the whole
representation to be unramified there.  Hsieh's enlarged-level
support supplies no such assertion; no claim is made that
every resulting component has nonzero monodromy either.
AC-G64-weight gives a separate exact obstruction to the
literal symmetric seed type.
The source \texttt{IsSeedAlt.above} and
\path{DenseThresholdCases/PrConc.lean},
\texttt{SeedPtOfAlt}, retain that symmetric assignment.
Finally the \texttt{IdealSpecializationPoint} and
\texttt{P1core}/\texttt{P3core} support requirements in
AC-G61 are not consequences of the new component statement.

\paragraph{Geometric dictionary.}
The faithful trace quotient distinguishes an actual simple
generic local system from an upper triangular extension with
the same trace.  The finite-cover and restriction language
has only the previously audited SGA~1, Expos\'e~V,
\S\S4 and~7 meaning.  The Hodge--Tate label obstruction is
a statement about the actual ordinary local filtration;
no geometric motive or new archival theorem is inferred
from the existence of Hecke support.
